\documentclass{situationist-preprint}

\title{The Location of Moral Facts}
\msid{MTE-01-c}
\msfamily{MTE-01}
\msdiscipline{Metaethics}
\mstype{Consequence}
\mskeywords{location problem; moral naturalism; idealized reactions; dispositions; reasons}

\begin{document}
\maketitle

\begin{abstract}
\citet{Jackson1998} poses the location problem for ethics: if the world is as the natural sciences describe it, where in that world are the moral facts? We derive the answer from the conclusion of MTE-01: moral facts are natural facts about idealized reactions. They are located among the counterfactual facts about how agents of a given nature would react under idealized conditions, and these have their categorical basis in the natural features of acts and in the psychology of the agents concerned. Moral facts are therefore located in two places at once, in the acts judged and in the judges, and nowhere outside the natural world.
\end{abstract}

\section{The derivation}\label{s:der}

By Theorem~4.2 of MTE-01-a, wrongness is idealized disapproval: an act is wrong just in case the reactions of agents who are fully informed, imaginatively engaged, free of errors of reasoning and impartial would converge on disapproval of it (Definitions~3.3 and 3.5 of MTE-01-a). A fact of this form is counterfactual. Counterfactual facts about reactions are located where dispositions are: in the categorical features that ground them \citep{Lewis1989}. Two kinds of feature ground the disposition to idealized disapproval of an act.

\begin{proposition}[Location]\label{p:loc}
Every moral fact is fixed by (i) the natural features of the act and its circumstances, including its effects on everyone it concerns, together with (ii) the natural features of the psychology of the agents whose idealized reactions are in question.
\end{proposition}

\begin{proof}
The idealized reaction to an act is the reaction its agents would have to complete natural information about it (Definition~3.3 of MTE-01-a). What the information is depends on (i); how agents would react to it depends on (ii). By the Naturalness assumption (Assumption~3.7 of MTE-01-a), idealized reactions are natural facts, so (ii) consists of natural facts; and by Theorem~6.2 of MTE-01-a, moral facts supervene on the natural facts. A moral fact therefore cannot vary while (i) and (ii) are held fixed.
\end{proof}

Clause (ii) concerns the nature of the agents, not their attitudes. MTE-02 argues that the content of moral requirements is relative to the nature of the agents concerned, and that every rational agent judging under idealized conditions reaches the same verdict; so the psychology in (ii) is that of a kind, not of an individual or a community. This is why moral facts are attitude-independent (Corollary~6.3 of MTE-01-a) although they are located partly in the psychology of judges: the actual attitudes of any agent may depart from the idealized ones.

\section{Two consequences}\label{s:cons}

First, no further location is needed. Non-naturalists hold that moral facts are sui generis and occupy a domain of their own \citep{ShaferLandau2003, Enoch2011}. By Proposition~\ref{p:loc}, once the natural facts in (i) and (ii) are fixed, nothing in such a domain could make a difference to what is wrong; a sui generis fact that made no difference would do no work.

Second, the location is the same whichever normative notion is taken as basic. MTE-11 holds that reasons are primitive and that all other normative notions are defined in terms of them. That is a claim about the order of definition, and it is compatible with the claim that what the normative concepts pick out is identical, a posteriori, with natural facts, just as the order of definition among chemical concepts is compatible with their reference to physical kinds. The facts that are reasons are located where Proposition~\ref{p:loc} locates the facts that are wrongs.

\begin{objection}
Counterfactuals about idealized agents have no categorical basis, since no actual agent is idealized.
\end{objection}

\begin{reply}
The counterfactual concerns how actual agents of a kind would react if their information, imagination, reasoning and impartiality were improved. Its basis is the sensibility they actually have, which conditions such as ignorance and partiality mask, as a fragile glass wrapped in padding remains fragile. The idealization specifies the conditions; the sensibility that would respond under them is actual.
\end{reply}

\section{Conclusion}

The location problem has a definite answer for the realism of MTE-01. Moral facts are located in the natural features of the acts judged and in the natural sensibility of the kind of agent who judges, and in nothing else.

\begin{thebibliography}{99}\small
\bibitem[Enoch(2011)]{Enoch2011} Enoch, D. (2011). \emph{Taking Morality Seriously: A Defense of Robust Realism}. Oxford University Press.
\bibitem[Jackson(1998)]{Jackson1998} Jackson, F. (1998). \emph{From Metaphysics to Ethics: A Defence of Conceptual Analysis}. Oxford University Press.
\bibitem[Lewis(1989)]{Lewis1989} Lewis, D. (1989). Dispositional theories of value. \emph{Proceedings of the Aristotelian Society, Supplementary Volume}, 63, 113--137.
\bibitem[Shafer-Landau(2003)]{ShaferLandau2003} Shafer-Landau, R. (2003). \emph{Moral Realism: A Defence}. Oxford University Press.
\end{thebibliography}

\end{document}
